\section{Постановка задачи}

\subsection{Входные данные и искомое решение}

Сеть представляется графом $G(t)=(V(t),E(t))$. Вершины соответствуют узлам, способным передавать и ретранслировать сообщения. Набор доступных каналов зависит от размещения устройств, параметров радиосвязи и состояния узлов. В имитационной модели координаты фиксированы в пределах прогона, а динамика задаётся отказами и восстановлением.

Входными данными являются координаты узлов, параметры LoRa, модель потерь, расписание прикладных сообщений и событий доступности. Каждое сообщение задаётся источником $s$, получателем $d$, временем создания и длиной полезной нагрузки. Общий канал полудуплексный: узел не может одновременно передавать и принимать.

Для алгоритмического уровня требуется найти путь $P$ между $s$ и $d$ с минимальной суммой неотрицательных весов:
\begin{equation}
  P^*=\arg\min_{P:s\leadsto d}\sum_{(u,v)\in P}w_{uv}.
  \label{eq:path}
\end{equation}
Наряду со стоимостью результата измеряются время вычисления, число посещённых вершин и число обработанных рёбер. Если пути нет, это фиксируется как отдельный исход. BFS минимизирует число переходов и используется как базовый метод, который не учитывает величину весов.

Для сетевого уровня требуется сопоставить стратегии передачи по доле доставленных сообщений, задержке, затратам эфира и энергии. Основной вопрос исследования: в каких условиях различия в доставке связаны с обслуживанием маршрутов и использованием радиоканала, даже если задача поиска пути на готовом графе решается эффективно?

\subsection{Метрики}

Коэффициент доставки определяется по уникальным прикладным сообщениям:
\begin{equation}
  \mathrm{PDR}=\frac{N_{\mathrm{del}}}{N_{\mathrm{gen}}},
\end{equation}
где $N_{\mathrm{gen}}$ --- число созданных сообщений, а $N_{\mathrm{del}}$ --- число сообщений, принятых конечным получателем. Повторные копии не увеличивают числитель. Задержка считается только для доставленных сообщений от создания до первого приёма адресатом. Поэтому задержку необходимо анализировать совместно с PDR: низкая задержка у небольшой доли успешных сообщений не означает высокое качество связи.

Служебные расходы характеризуются долей суммарного времени передач:
\begin{equation}
  H=\frac{T_{\mathrm{ctrl}}}{T_{\mathrm{ctrl}}+T_{\mathrm{data}}}.
\end{equation}
Это отношение сумм длительностей кадров, а не доля неперекрывающегося физического времени наблюдения. Кроме $H$ учитывается число передач на одно доставленное сообщение. У лавинной схемы может отсутствовать отдельный служебный трафик маршрутизации, но повторная передача данных всё равно расходует эфир.

Энергоэффективность задаётся как
\begin{equation}
  E_{\mathrm{bit}}=\frac{E_{\mathrm{net}}}{8L N_{\mathrm{del}}},
\end{equation}
где $E_{\mathrm{net}}$ --- суммарная энергия узлов, а $L$ --- длина полезной нагрузки в байтах. При отсутствии доставленных сообщений показатель не определён. Дополнительно рассматриваются полезная скорость доставки, время до истощения первой батареи и изменение PDR при отказах.

Окончательный выбор стратегии предполагает компромисс между надёжностью, задержкой и энергией. Единственная взвешенная сумма этих показателей не вводится: предпочтение зависит от требований конкретного сценария оффлайн-связи.
